\chapter{Project Management and Development Evidence}
\section{Work Completed}
The isolated copy preserves original history/documents. Completed work includes complete feature/family evidence, independent prose/JSON checks, strict mode, rejection/fallback reporting, local LLM transport/evaluation plumbing, both live interfaces, restored viewer runtime, authoritative-record validation, an explicit prototype profile and a separate retrospective risk comparator. The selected suite has 212 passes and no skips; local evidence covers 372 records, 62 live HTTP cases and 228 comparator arithmetic probes. CI jobs are prepared.
\section{Work Remaining}
The real fresh Alpaca scientific run and original-publisher paired social comparison now complete, with retained raw execution, source/input identity and replay evidence. Cardiff leads the small AI-reference benchmark; final adequacy criteria, independent validation and deployment selection remain open. Source/raw labels, exclusions and the original V3 chronology are preserved.

Current gaps are institutional report/deck inspection and completion, genuine member contributions and review/rehearsal, current repository/remote CI evidence, live LLM usefulness and agreed risk semantics. The original September bytes and exact scientific lock are not reproduced. Student inference/training and broader acceptance work remain later milestones.

\section{Contribution of Each Team Member}
Section~\ref{sec:planned-work-division} contains the proposal's planned member assignments. Ahmed covers the point-in-time harness, leakage/timestamp/parity work and deployment; Rafay covers risk profiling, sentiment, degraded operation, acceptance and frontend; Soban covers the decision engine, calibration, SHAP, verifier and later QLoRA. Each member must record actual implementation, review, testing and defence preparation. The planned table does not establish completion or authorship of assistant-supported work.
\section{Repository and Version Control Evidence}
The consolidated repository is on \texttt{main}. Original baseline \texttt{1dc24e0} and its isolated branch are preserved in history; local commit \texttt{93d0d60} records the fresh scientific increment. Current verification evidence records source/data hashes. Upstream access remains unverified; local workflow configuration is not evidence of a successful remote CI run.
\section{Continuous Integration and Delivery Status}
The GitHub Actions configuration runs two Python 3.11 jobs on Windows and Ubuntu after a push, pull request or manual dispatch. Both call the same versioned \texttt{run\_mid1\_ci.py} entry point used for local replay. The portable job executes standard-library regressions, saved-record validation, preview generation, a plan-only explainer check, comparator arithmetic and a real localhost HTTP exercise. The UI job runs scoped source checks and the legacy, presenter and recorded-study selector suites. Each retains logs and a byte-bound source/artifact receipt; CI artifacts are configured for fourteen-day retention.

The Windows local replay passes both jobs: 165 portable tests and seven CI source-ownership/metadata guards, 62 live HTTP checks, 372 validated historical records and 228 comparator probes, plus 42 UI tests in separate overlapping scopes. Temporary component metadata is restored after each job. The replay uses the already prepared local runtime; it does not establish a fresh dependency installation or Ubuntu/GitHub execution. The consolidated repository still has no configured remote. Actual remote CI remains open, and no continuous deployment, hosting target, model training or paid inference is executed by this workflow.
\section{Milestones and Commit Evidence}
The checkpoint distinguishes the initial context increment from verifier/local-demo work. Historical experiments are preserved separately from current test/HTTP evidence. Reproduction instructions appear in the repository appendix.
\section{Tools and Tool-Assisted Development}
The coding assistant contributed code, regression tests, drafts and programmatic validation. Earlier handoffs declare GPT-6.1 Sol/Xhigh as selected session configuration; the exact serving revision is not provider-authenticated. Student review, source verification, institutional declarations and accurate tool disclosure remain necessary. Student authorship is not inferred from assistant work.
\section{External Components and Dependencies}
The baseline uses FinBERT, XGBoost, isotonic calibration, TreeSHAP and Streamlit. The portable demo uses standard Python HTTP/JSON facilities and existing record/verifier code. The retained social labels were generated by the assistant in one actual session pass and disclosed as AI references. Six recorded assistant-teacher explanation completions are retained separately. No API/local student inference or human annotation is claimed.
\section{Risks and Mitigation}
Both interfaces expose recorded evidence and failure cases. The restored scientific runtime executes a separate fresh snapshot; original-byte and exact-lock limitations remain disclosed. Unknown deadline/owners affect scheduling; unavailable model access prevents measuring LLM usefulness. Constant-Buy classifier output prevents a predictive-skill claim. Prototype profile-to-blend assumptions require review and do not establish suitability.
\section{Security and Privacy Considerations}
The local demo serves historical decisions on loopback and has no trading/account access. Profile persistence, production authentication and deployment security are beyond this increment.
\section{Current Limitations}
Physical projector review, remote GitHub CI, real LLM evaluation, personalised classifier modelling and deployment remain pending. The current browser journey and compiled report/deck layouts have separate review receipts. New science and the original social-classifier comparison are recorded separately. The risk comparator remains an unvalidated prototype. Institutional metadata, signatures, genuine team review and remaining research are incomplete.

\section{Teacher Pilot Evidence}
The six recorded examples, four-training/two-validation split, serving-identity limitations and untrained Phi-3-mini student are documented once in Section~\ref{sec:teacher-pilot}. No training efficacy or human-comprehension result is inferred.
